\documentclass{beamer}

\usepackage{cuhkbeamer}

\begin{document}

% ====================== Slide 1: TITLE SLIDE ======================
\begin{frame}
  \title{Demand Modeling in Urban Air Mobility}
  \author{Chiu Yu KO}
  \institute{Low Altitude Economics}
  \date{Week 4}
  \titlepage
\end{frame}
% ====================== Slide 2: FORMAT & OBJECTIVES ======================
\begin{frame}[t]
\frametitle{Learning Objectives}
By the end of this class, you should be able to:
\begin{itemize}
\item \textbf{Build} a simple demand forecast from potential trips, corridor share, mode share, and operating constraints.
\item \textbf{Estimate} willingness to pay using generalized cost and interpret a simple demand curve.
\item \textbf{Explain} how own-price elasticity and customer heterogeneity affect pricing decisions.
\item \textbf{Incorporate} public acceptance as a transparent scenario assumption rather than a fixed universal multiplier.
\item \textbf{Stress-test} demand and revenue under weather, capacity, and policy scenarios.
\end{itemize}
\end{frame}

% ====================== Slide 3: ICEBREAKER ======================
\begin{frame}[t]
\frametitle{Would Customers Actually Choose It?}
\textbf{Imagine a proposed airport air-mobility service in Hong Kong.}
\begin{itemize}
\item The aircraft may be technically capable of completing the trip quickly.
\item Customers still compare the complete journey with rail, taxi, bus, and private car alternatives.
\end{itemize}
\vspace{0.3cm}
The commercial question is not simply whether the service can fly. It is whether enough customers will choose it at a viable price.
\begin{itemize}
\item \textbf{Quick poll (2 min):} Which factor is likely to matter most: price, door-to-door time, reliability, convenience, or public confidence?
\end{itemize}
\end{frame}

% ====================== Slide 4: CONCEPT PRIMER 1 ======================
\begin{frame}[t]
\frametitle{Concept Primer: WTP and Demand Curves}
\textbf{Willingness to pay (WTP)} is the maximum amount a particular customer would pay for a service rather than go without it or choose an alternative.
\begin{itemize}
\item Customers have different WTP because they value time, reliability, comfort, and urgency differently.
\item A customer does not necessarily disappear from the entire market when price exceeds WTP; they may choose another mode.
\end{itemize}
\textbf{A demand curve} shows the quantity customers are expected to purchase at different prices, holding other important conditions fixed.
\begin{itemize}
\item It helps managers compare volume, revenue, capacity requirements, and contribution at alternative prices.
\end{itemize}
\end{frame}


% ====================== NEW SLIDE: APPLIED TOOL - WTP & DEMAND CURVE ======================
\begin{frame}[t]
\frametitle{Applied Tool: A Simple Linear Demand Curve}
\begin{block}{The Formula}
\[Q=a-bP\]
\end{block}
\begin{itemize}
\item $Q$ = expected quantity demanded during a defined period
\item $P$ = price per unit
\item $a$ = estimated demand at a zero price under the model
\item $b$ = change in quantity associated with a HK\$1 price increase
\item The choke price is $P=a/b$, where estimated demand reaches zero.
\end{itemize}
\textbf{Managerial caution:} A survey does not automatically produce a reliable demand curve. Managers should use multiple observations, realistic alternatives, and validation against actual behaviour where possible.
\end{frame}

% ====================== NEW SLIDE: BACK-OF-THE-ENVELOPE (WTP \& DEMAND CURVE) ======================
\begin{frame}[t,shrink=10]
\frametitle{Back-of-the-Envelope: Comparing Prices for an Airport Service}
\textbf{Illustrative demand model:} $Q=8{,}000-8P$. Demand falls by 800 passenger-trips for each HK\$100 price increase.
\begin{columns}[T]
\begin{column}{0.48\textwidth}\begin{block}{Price = HK\$300}
\begin{itemize}
\item $Q=8{,}000-8(300)=5{,}600$
\item Revenue $=300\times5{,}600$
\item \textbf{Revenue = HK\$1.68 million}
\end{itemize}\end{block}\end{column}
\begin{column}{0.48\textwidth}\begin{block}{Price = HK\$500}
\begin{itemize}
\item $Q=8{,}000-8(500)=4{,}000$
\item Revenue $=500\times4{,}000$
\item \textbf{Revenue = HK\$2.00 million}
\end{itemize}\end{block}\end{column}
\end{columns}
\textit{Takeaway: HK\$500 maximizes revenue for this particular linear model. A profit-maximizing price would also require cost and capacity information.}
\vfill{\tiny Illustrative assumptions for classroom analysis; not a survey result or market forecast.}
\end{frame}


% ====================== Slide 5: CONCEPT PRIMER 2 (REVISED) ======================
\begin{frame}[t]
\frametitle{Concept Primer: What is Demand Modelling?}
Demand modelling uses evidence and explicit assumptions to estimate how many customers may choose a service under different conditions.
\begin{itemize}
\item It is not a precise prediction of exactly how many people will travel.
\item It combines market size, trip patterns, competing modes, price, time, reliability, access, and customer preferences.
\item The forecast should be presented as a range or set of scenarios rather than a single certain number.
\end{itemize}
\textbf{Link to Week 3:} Does forecast demand fit the price, capacity, and activity required to break even?
\end{frame}

% ====================== Slide 6: NASA FRAMEWORK ======================
\begin{frame}[t]
\frametitle{A Four-Step Travel-Demand Framework}
A traditional transport-planning framework can organize an LAE demand forecast:
\begin{enumerate}
\item \textbf{Trip generation:} How many relevant trips may occur?
\item \textbf{Trip distribution:} Which origins and destinations do those trips connect?
\item \textbf{Mode choice:} What share may choose LAE rather than rail, taxi, bus, or car?
\item \textbf{Network assignment:} How would the chosen trips use the available network, routes, and facilities?
\end{enumerate}
\textit{NASA-sponsored UAM studies use detailed demand and market models, but the four-step structure is a general transport-planning framework rather than a unique ``NASA formula.''}
\end{frame}


% ====================== Slide 7: APPLIED TOOL - NASA 4-STEP DEMAND FORECAST ======================
\begin{frame}[t]
\frametitle{Applied Tool: A Simple Demand-Funnel Approximation}
\begin{block}{A Managerial Screening Formula}
\[
\text{Potential LAE Passenger-Trips}
=\text{Relevant Trips}\times\text{Corridor Share}\times\text{Mode Share}
\]
\end{block}
\begin{block}{Capacity Check}
\[
\text{Served Passenger-Trips}
=\min\{\text{Potential Demand},\text{Available Capacity}\}
\]
\end{block}
\begin{itemize}
\item The multiplication is a simplified demand funnel, not the full four-step transport model.
\item Route assignment allocates trips across a network; it should not normally be treated as an arbitrary demand-reduction percentage.
\item Keep passenger-trips, flights, and seats separate throughout the model.
\end{itemize}
\end{frame}


% ====================== Slide 8: BACK-OF-THE-ENVELOPE (NASA 4-STEP) ======================
\begin{frame}[t,shrink=12]
\frametitle{Back-of-the-Envelope: Demand and Capacity for an Airport Corridor}
\textbf{Assumptions:} 20,000 relevant daily trips; 60\% involve the selected corridor; estimated LAE mode share = 35\%.
\begin{columns}[T]
\begin{column}{0.48\textwidth}\begin{block}{Potential Demand}
\begin{itemize}
\item $20{,}000\times0.60\times0.35$
\item \textbf{4,200 passenger-trips/day}
\item This is customer demand before checking operating capacity.
\end{itemize}\end{block}\end{column}
\begin{column}{0.48\textwidth}\begin{block}{Capacity-Constrained Service}
\begin{itemize}
\item Assume 300 flights/day
\item Four seats and 80\% target load factor
\item Capacity $=300\times4\times0.80$
\item \textbf{960 passenger-trips/day}
\end{itemize}\end{block}\end{column}
\end{columns}
\textit{Takeaway: Potential demand may exceed available capacity. The business model must distinguish customers who want to travel from customers the system can actually serve.}
\vfill{\tiny Illustrative assumptions for classroom analysis; not current corridor data.}
\end{frame}

% ====================== Slide 7: CONCEPT PRIMER 3 ======================
\begin{frame}[t]
\frametitle{Concept Primer: Generalized Cost}
Generalized cost combines the main burdens of a complete journey into a common monetary measure.
\begin{itemize}
\item fare and related charges;
\item access, waiting, boarding, in-vehicle, and transfer time;
\item reliability, inconvenience, and other service attributes where these can be estimated credibly.
\end{itemize}
\textbf{Managerial use:} Compare door-to-door alternatives on a consistent basis.
\vspace{0.25cm}
\textit{A lower generalized cost increases the attractiveness of a mode, but does not imply that every customer will choose it. Preferences differ and forecasts remain probabilistic.}
\end{frame}

% ====================== Slide 8: STANDALONE FORMULA ======================
\begin{frame}[t]
\frametitle{Applied Tool: Generalized Cost}
\begin{block}{A Simple Formula}
\[
GC=\text{Fare}+VOT\times\text{Door-to-Door Time}+\text{Other Monetized Burdens}
\]
\end{block}
\begin{itemize}
\item $VOT$ converts travel time into money for comparison.
\item Other burdens may include unreliability, difficult access, transfers, or screening, provided the values are estimated transparently.
\item Avoid double counting. A weather delay should not appear in both expected travel time and a separate penalty unless the terms capture different effects.
\item Customers are more likely, not certain, to select alternatives with lower generalized cost.
\end{itemize}
\end{frame}

% ====================== Slide 9: BACK-OF-ENVELOPE (GENERALIZED COST) ======================
\begin{frame}[t,shrink=10]
\frametitle{Back-of-the-Envelope: Comparing Door-to-Door Generalized Cost}
\textbf{Assumption:} VOT = HK\$15 per minute. Other monetized burdens are illustrative.
\begin{columns}[T]
\begin{column}{0.48\textwidth}\begin{block}{Ground Taxi}
\begin{itemize}
\item Fare: HK\$350
\item Door-to-door time: 45 minutes
\item Other burden: HK\$50
\item $GC=350+15(45)+50$
\item \textbf{GC = HK\$1,075}
\end{itemize}\end{block}\end{column}
\begin{column}{0.48\textwidth}\begin{block}{Air-Mobility Service}
\begin{itemize}
\item Fare: HK\$800
\item Door-to-door time: 12 minutes
\item Other burden: HK\$80
\item $GC=800+15(12)+80$
\item \textbf{GC = HK\$1,060}
\end{itemize}\end{block}\end{column}
\end{columns}
\textit{Takeaway: The air service is slightly more attractive under these assumptions. The difference is small, so changes in waiting time, reliability, or perceived risk could reverse the result.}
\vfill{\tiny Illustrative assumptions for classroom analysis.}
\end{frame}

% ====================== Slide 10: SESSION 2 TITLE ======================
\begin{frame}[c]
  \begin{center}
    \huge\color{cuhkpurple}\textbf{Session 2}\\
    \vspace{0.5cm}
    \Large WTP Estimation, Elasticities \& Heterogeneity
  \end{center}
\end{frame}

% ====================== Slide 11: CONCEPT PRIMER 4 ======================
\begin{frame}[t]
  \frametitle{Concept Primer: "Elasticity of Demand"}
  
  \textbf{What happens when you change your ticket price?}
  \begin{itemize}
      \item If you raise the price of a flying taxi by 10\%, will you lose 5\% of your customers, or 50\% of your customers?
      \item \textbf{Elasticity} measures this exact sensitivity.
  \end{itemize}

  \vspace{0.3cm}
  \textbf{Two Types of Elasticity:}
  \begin{itemize}
      \item \textbf{Own-Price Elasticity:} How your demand changes when \textit{you} change your price.
      \item \textbf{Cross-Elasticity:} How your demand changes when \textit{your competitor} changes their price. (e.g., If the MTR raises its fare by HK\$10, how many people suddenly switch to your flying taxi?)
  \end{itemize}
\end{frame}

% ====================== Slide 12: STANDALONE FORMULA ======================
\begin{frame}[t]
\frametitle{Applied Tool: Price Elasticity of Demand}
\begin{block}{The Formula}
\[
E_d=\frac{\%\text{ change in quantity demanded}}{\%\text{ change in price}}
\]
\end{block}
\begin{itemize}
\item Own-price elasticity is usually negative because demand tends to fall as price rises.
\item Managers often discuss its absolute value: $|E_d|<1$ is inelastic and $|E_d|>1$ is elastic.
\item The elasticity may change across prices, customer groups, time periods, and competing alternatives.
\item Small-change approximations are less reliable for large price changes.
\end{itemize}
\textbf{Managerial caution:} Emergency or public-service demand does not imply unlimited ability to charge. Budgets, procurement rules, contracts, and public policy still constrain price.
\end{frame}

% ====================== Slide 13: BACK-OF-ENVELOPE (ELASTICITY) ======================
\begin{frame}[t,shrink=10]
\frametitle{Back-of-the-Envelope: How Elasticity Affects Revenue}
\textbf{Starting point:} Price = HK\$500; volume = 100 passenger-trips; revenue = HK\$50,000. Consider a 20\% price reduction.
\begin{columns}[T]
\begin{column}{0.48\textwidth}\begin{block}{Approximate $|E_d|=0.5$}
\begin{itemize}
\item Approximate volume increase: 10\%
\item New volume: 110
\item New revenue: $400\times110$
\item \textbf{HK\$44,000}
\end{itemize}\end{block}\end{column}
\begin{column}{0.48\textwidth}\begin{block}{Approximate $|E_d|=2.0$}
\begin{itemize}
\item Approximate volume increase: 40\%
\item New volume: 140
\item New revenue: $400\times140$
\item \textbf{HK\$56,000}
\end{itemize}\end{block}\end{column}
\end{columns}
\textit{Takeaway: A discount can raise or lower revenue depending on demand response. Profit also depends on capacity and incremental cost.}
\vfill{\tiny Illustrative constant-elasticity approximation applied to a relatively large price change.}
\end{frame}

% ====================== Slide 14: DEMAND HETEROGENEITY ======================
\begin{frame}[t]
\frametitle{Demand Heterogeneity}
Customers differ in urgency, value of time, budget, reliability needs, accessibility requirements, booking flexibility, and willingness to adopt a new service.
\begin{itemize}
\item A single corridor may therefore contain several customer segments with different demand curves.
\item Managers may offer different products, booking conditions, service levels, memberships, or time-based prices.
\item Segmentation must be operationally feasible, transparent, compatible with applicable law and policy, and acceptable to customers.
\end{itemize}
\textbf{Managerial objective:} Match the product and price to different needs without assuming that one label perfectly predicts behaviour.
\end{frame}


% ====================== NEW SLIDE: APPLIED TOOL - DEMAND HETEROGENEITY & PRICE DISCRIMINATION ======================
\begin{frame}[t]
\frametitle{Applied Tool: Revenue by Customer Segment}
\begin{block}{The Formula}
\[
\text{Total Revenue}=\sum_{i=1}^{n}P_iQ_i(P_i)
\]
\end{block}
\begin{itemize}
\item $P_i$ is the price or product offer for segment $i$.
\item $Q_i(P_i)$ is that segment's expected demand at the offered price.
\item Different offers may reflect booking flexibility, priority, refundability, service level, timing, or membership rather than personal characteristics.
\item Revenue maximization should be checked against cost, capacity, fairness, transparency, and legal constraints.
\end{itemize}
\end{frame}

% ====================== NEW SLIDE: BACK-OF-THE-ENVELOPE (DEMAND HETEROGENEITY) ======================
\begin{frame}[t,shrink=10]
\frametitle{Back-of-the-Envelope: Two Product Offers}
\textbf{Scenario:} Compare a uniform fare with two differentiated products for an airport corridor.
\begin{columns}[T]
\begin{column}{0.48\textwidth}\begin{block}{Uniform Fare: HK\$500}
\begin{itemize}
\item Flexible-business demand: 4,000
\item Advance-purchase demand: 2,000
\item \textbf{Revenue = HK\$3.00 million}
\end{itemize}\end{block}\end{column}
\begin{column}{0.48\textwidth}\begin{block}{Differentiated Products}
\begin{itemize}
\item Flexible fare: HK\$800, demand 3,800
\item Restricted advance fare: HK\$300, demand 3,500
\item \textbf{Revenue = HK\$4.09 million}
\end{itemize}\end{block}\end{column}
\end{columns}
\textit{Takeaway: Product differentiation can increase revenue under the assumed demand responses. Costs, capacity, cannibalization, fairness, and implementation must also be assessed.}
\vfill{\tiny Hypothetical figures for classroom analysis; not a recommendation for pricing based on protected personal characteristics.}
\end{frame}

% ====================== Slide 15: WHAT IS AiRMOUR? ======================
\begin{frame}[t]
\frametitle{What is AiRMOUR?}
AiRMOUR was an EU-funded research and innovation project supporting sustainable urban air mobility, with a strong focus on emergency medical services.
\begin{itemize}
\item Its guidebook brings together urban design and mobility, aviation safety, public acceptance, and integration-process management.
\item It treats public and stakeholder acceptance as part of planning and implementation, not merely as a numerical demand discount.
\item Its evidence can help managers identify questions, engagement needs, risks, and scenarios for local analysis.
\end{itemize}
\vfill{\tiny Source: AiRMOUR, \textit{Guidebook for Urban Air Mobility Integration} (2023).}
\end{frame}

% ====================== Slide 16: AiRMOUR SEGMENTATION (HK FOCUS) ======================
\begin{frame}[t]
\frametitle{Public Acceptance: A Hong Kong Application}
Acceptance may differ by mission, location, affected group, perceived benefit, noise, privacy, safety, fairness, and trust.
\begin{itemize}
\item Emergency medical use may be viewed differently from leisure, passenger, inspection, or commercial delivery use.
\item However, ``medical'' does not imply full support, perfectly inelastic demand, or faster approval.
\item Non-medical applications should not be assigned one universal level of opposition.
\item Engagement should include users, nearby residents, public agencies, existing transport providers, and other affected stakeholders.
\end{itemize}
\textbf{Managerial lesson:} Estimate acceptance locally and use scenarios until reliable evidence becomes available.
\end{frame}


% ====================== NEW SLIDE: APPLIED TOOL - AiRMOUR ACCEPTANCE ADJUSTMENT ======================
\begin{frame}[t]
\frametitle{Applied Tool: Acceptance Scenarios}
\begin{block}{A Transparent Scenario Adjustment}
\[
\text{Scenario Demand}=\text{Base Demand}\times\text{Acceptance Adjustment}
\]
\end{block}
\begin{itemize}
\item The adjustment is an analyst-defined scenario input, not an official AiRMOUR coefficient.
\item Use several values, such as 0.6, 0.8, and 1.0, and explain the evidence or judgment behind them.
\item Avoid double counting if public acceptance already affects stated mode share, route availability, or generalized cost.
\item Over time, replace assumed adjustments with survey, pilot, complaint, booking, and revealed-behaviour evidence.
\end{itemize}
\end{frame}

% ====================== NEW SLIDE: BACK-OF-THE-ENVELOPE (AiRMOUR) ======================
\begin{frame}[t,shrink=10]
\frametitle{Back-of-the-Envelope: Acceptance Scenario Analysis}
\textbf{Base forecast:} 4,200 potential passenger-trips per day before an acceptance adjustment.
\begin{columns}[T]
\begin{column}{0.48\textwidth}\begin{block}{Higher-Acceptance Scenario}
\begin{itemize}
\item Adjustment: 0.9
\item Demand: $4{,}200\times0.9$
\item \textbf{3,780 passenger-trips/day}
\end{itemize}\end{block}\end{column}
\begin{column}{0.48\textwidth}\begin{block}{Lower-Acceptance Scenario}
\begin{itemize}
\item Adjustment: 0.6
\item Demand: $4{,}200\times0.6$
\item \textbf{2,520 passenger-trips/day}
\end{itemize}\end{block}\end{column}
\end{columns}
\textit{Takeaway: Acceptance assumptions can materially change forecasts. They should be visible, justified, varied, and not presented as measured facts.}
\vfill{\tiny Illustrative scenario factors; not coefficients published by AiRMOUR.}
\end{frame}

% ====================== Slide 16: SESSION 3 TITLE ======================
\begin{frame}[c]
  \begin{center}
    \huge\color{cuhkpurple}\textbf{Session 3}\\
    \vspace{0.5cm}
    \Large Scenario Analysis, Uncertainty \& Shocks
  \end{center}
\end{frame}

% ====================== Slide 17: CONCEPT PRIMER 5 ======================
\begin{frame}[t]
\frametitle{Concept Primer: Scenario Analysis}
A demand forecast should examine plausible changes rather than assume one perfectly stable future.
\begin{itemize}
\item \textbf{Demand shock:} Fewer customers travel or choose the service.
\item \textbf{Capacity shock:} Weather, maintenance, staffing, or operating restrictions reduce service availability.
\item \textbf{Price or cost shock:} Fares, competing-mode prices, energy, insurance, or site costs change.
\item \textbf{Policy or acceptance change:} New conditions alter routes, operating hours, or customer confidence.
\end{itemize}
\textbf{Managerial purpose:} Identify which assumptions matter most and prepare responses before committing large investments.
\end{frame}

% ====================== Slide 18: STANDALONE FORMULA ======================
\begin{frame}[t]
\frametitle{Applied Tool: Revenue Under Demand and Capacity Shocks}
\begin{block}{Served Quantity}
\[
Q_{served}=\min\{Q_{demand},Q_{capacity}\}
\]
\end{block}
\begin{block}{Revenue}
\[
R=P\times Q_{served}+\text{Committed Revenue}-\text{Refunds and Credits}
\]
\end{block}
\begin{itemize}
\item A weather event may reduce capacity rather than customer demand, although repeated unreliability can later affect demand.
\item A policy change may affect price, capacity, cost, or demand through different channels.
\item Model each channel explicitly rather than applying one unexplained ``shock penalty'' to revenue.
\end{itemize}
\end{frame}

% ====================== Slide 19: BACK-OF-ENVELOPE (SHOCKS) ======================
\begin{frame}[t,shrink=10]
\frametitle{Back-of-the-Envelope: A Weekly Capacity Shock}
\textbf{Assumptions:} Demand and normal capacity = 100 trips/week; price = HK\$500; fixed cost = HK\$40,000/week. Variable cost is omitted for simplicity.
\begin{columns}[T]
\begin{column}{0.48\textwidth}\begin{block}{Normal Week}
\begin{itemize}
\item Served trips: 100
\item Revenue: HK\$50,000
\item \textbf{Margin before variable cost: HK\$10,000}
\end{itemize}\end{block}\end{column}
\begin{column}{0.48\textwidth}\begin{block}{30\% Capacity Loss}
\begin{itemize}
\item Served trips: 70
\item Revenue: HK\$35,000
\item \textbf{Margin before variable cost: $-$HK\$5,000}
\end{itemize}\end{block}\end{column}
\end{columns}
\textit{Takeaway: Possible responses include reserves, flexible capacity, alternative transport, committed revenue, insurance, or revised pricing. Raising the fare to HK\$600 alone would produce HK\$42,000 only if all 70 customers still travelled.}
\vfill{\tiny Illustrative assumptions; excludes variable cost, refunds, service credits, and demand response to a price change.}
\end{frame}

% ====================== Slide 20: HK VS GBA COMPARISON (REVISED) ======================
\begin{frame}[t]
\frametitle{Hong Kong and the GBA: Different Demand Opportunities}
\begin{center}\small
\begin{tabular}{l|cc}
\textbf{Dimension} & \textbf{Hong Kong possibility} & \textbf{Wider GBA possibility}\\
\hline
Trip patterns & Dense, complex, constrained & Diverse urban and regional\\
Customer value & Time-sensitive, specialized & Larger and more varied\\
Alternatives & Strong public transport & Varies substantially by city\\
Scale & Limited local geography & Broader network potential\\
Key question & Value per trip & Volume and network reach\\
\end{tabular}
\end{center}
\vspace{0.3cm}
\textit{These are hypotheses for analysis, not universal conclusions. Forecasts should use route-specific data on customers, alternatives, access, capacity, and regulation.}
\end{frame}


% ====================== Slide 21: EXERCISE INSTRUCTIONS ======================
\begin{frame}[t]
  \frametitle{In-Class Exercise: Interactive Pricing Simulation}
  
  \textbf{Activity: The Pricing Sweet Spot}
  \begin{itemize}
    \item \textbf{Format:} Groups of 4–5 (Laptops required)
    \item \textbf{Scenario:} You are managing a HK Airport Shuttle corridor.
    \item \textbf{Task:} Using the provided Excel template:
    \begin{enumerate}
        \item Set prices for Segment A (Business) and Segment B (Tourists).
        \item Observe how the Elasticity multiplier changes your total volume.
        \item Optimise your prices to achieve maximum total revenue.
    \end{enumerate}
    \item \textbf{Timing:} 15 minutes hands-on + 5 minute share-out.
  \end{itemize}
\end{frame}

% ====================== Slide 22: EXERCISE TEMPLATE ======================
\begin{frame}[t]
  \frametitle{Exercise Template (Visual)}
  
  \begin{center}
      \textit{[Distribute digital Excel Spreadsheet to student groups]}
      
      \vspace{0.5cm}
      
      \textbf{The Dashboard View:} \\
      \vspace{0.2cm}
      Input: Base Price | Executive Discount (\%) | Tourist Discount (\%) \\
      Output: Executive Volume | Tourist Volume | Total Revenue \\
      \vspace{0.5cm}
      \textit{Hint: Don't give discounts to executives with inelastic demand!}
  \end{center}
\end{frame}

% ====================== Slide 23: KEY TAKEAWAYS ======================
\begin{frame}[t]
\frametitle{Key Takeaways}
\begin{itemize}
\item Demand modelling produces evidence-based scenarios, not exact predictions.
\item Distinguish potential passenger demand, available capacity, flights, and passengers actually served.
\item Generalized cost compares complete journeys, while mode choice remains heterogeneous and probabilistic.
\item Elasticity helps assess price changes, but revenue and profit also depend on capacity and cost.
\item Customer segmentation can support differentiated products, subject to transparency, fairness, feasibility, and applicable rules.
\item Public acceptance and external shocks should be modelled through explicit, justified assumptions rather than universal multipliers.
\end{itemize}
\end{frame}

% ====================== Slide 24: ASSIGNMENT ======================
\begin{frame}[t]
\frametitle{Week 4 Assignment Briefing}
\textbf{Demand Forecast Memo and Model (maximum 1,200 words plus spreadsheet)}
\begin{itemize}
\item \textbf{Due:} Beginning of Week 5 (group assignment)
\item \textbf{Task:} Extend the previous business model with a transparent demand forecast.
\item \textbf{Requirements:}
\begin{enumerate}
\item Define the relevant trip market, corridor, time period, and customer groups.
\item Estimate potential demand and mode share using generalized cost or a clearly explained alternative.
\item Compare forecast demand with available seats and flights.
\item Estimate price sensitivity for two product or customer segments and state the limitations.
\item Run one acceptance scenario and one weather or capacity shock.
\item Recommend a pricing or capacity response and explain its risk.
\end{enumerate}
\item \textbf{Grading focus:} Economic reasoning (40\%), transparent calibration (30\%), scenario analysis (20\%), professionalism (10\%).
\end{itemize}
\end{frame}

% ====================== Slide 25: PREVIEW WEEK 5 ======================
\begin{frame}[t]
  \frametitle{Preview of Week 5}
  
  \begin{center}
      \Large \textbf{Infrastructure Economics and Management}
  \end{center}
  
  \vspace{1cm}
  \textbf{Preview question for next week:} \\
  \textit{How will these demand forecasts drive exactly where we place million-dollar vertiports in land-scarce Hong Kong?}
\end{frame}

% ====================== Slide 26: THANK YOU ======================
\begin{frame}[t]
  \frametitle{Thank You + Q\&A}
  
  \textbf{Pre-Class Readings Reminder (for Week 5):}
  \begin{itemize}
      \item \textit{McKinsey: Perspectives on AAM (Infrastructure sections)}
      \item \textit{NASA UAM Market Study (Site-selection sensitivities)}
  \end{itemize}
  
  \vspace{1cm}
  \begin{center}
      \textbf{Questions?} \\
      \small (Final 15 minutes reserved for extended Q\&A)
  \end{center}
\end{frame}

\end{document}